% Basic LaTeX layout for presentation of TFG defense
\documentclass[12pt,a4paper]{article}
\usepackage[utf8]{inputenc}
\usepackage[T1]{fontenc}
\usepackage[spanish]{babel}
\usepackage{lmodern}
\usepackage{geometry}
\geometry{margin=2.5cm}
\usepackage{hyperref}
\hypersetup{
	colorlinks=true,
	linkcolor=blue,
	urlcolor=blue,
	pdftitle={Defensa TFG},
	pdfauthor={Autor}
}

\title{Guion de la Defensa del TFG}
\author{Nombre del Estudiante}
\date{\today}

\begin{document}

\maketitle

\begin{abstract}
Breve resumen del contenido de la defensa: objetivos, metodología y resultados.
\end{abstract}

\section{Introducción}
Presentación del tema, motivación y planteamiento del problema.

\section{Objetivos}
\begin{itemize}
	\item Objetivo general.
	\item Objetivos específicos.
\end{itemize}

\section{Metodología}
Descripción de métodos, herramientas y etapas del trabajo.

\section{Resultados}
Resumen de los principales resultados obtenidos.

\section{Conclusiones}
Conclusiones y trabajo futuro.

\section*{Bibliografía}
% Use \bibliographystyle and \bibliography if .bib is available
\begin{thebibliography}{9}
\bibitem{ref1} Autor. \textit{Título}. Editorial, Año.
\end{thebibliography}

\end{document}
